\documentclass[11pt,a4paper]{article}
\usepackage[utf8]{inputenc}
\usepackage[T1]{fontenc}
\usepackage[italian]{babel}
\usepackage{lmodern}
\usepackage{microtype}
\usepackage[margin=2cm,headheight=14pt]{geometry}
\usepackage{amsmath,amssymb,amsthm,mathtools,amsfonts,bm,siunitx}
\usepackage{booktabs,tabularx,enumitem,float}
\usepackage{xcolor}
\usepackage[most]{tcolorbox}
\usepackage{fancyhdr}
\usepackage{listings}
\usepackage{hyperref}
\usepackage{bookmark}

\definecolor{navyblue}{RGB}{10, 45, 90}
\definecolor{coolblue}{RGB}{28, 110, 164}
\definecolor{forestgreen}{RGB}{20, 115, 60}
\definecolor{warmamber}{RGB}{195, 110, 10}
\definecolor{crimsonred}{RGB}{175, 30, 45}
\definecolor{subtlegray}{RGB}{245, 247, 250}
\definecolor{codebg}{RGB}{240, 242, 245}

\hypersetup{
    colorlinks=true,
    linkcolor=navyblue,
    urlcolor=coolblue,
    citecolor=forestgreen
}

\pagestyle{fancy}
\fancyhf{}
\fancyhead[L]{\textbf{\textsf{\textcolor{navyblue}{Statistica I --- Soluzione Ufficiale Dettagliata}}}}
\fancyhead[R]{\textsf{\textcolor{coolblue}{I Appello (08/06/2022)}}}
\fancyfoot[C]{\thepage}
\renewcommand{\headrulewidth}{0.6pt}

\newtcolorbox{problemabox}[1]{
    enhanced,
    breakable,
    colback=white,
    colframe=navyblue,
    coltitle=white,
    fonttitle=\bfseries\sffamily,
    title={#1},
    attach boxed title to top left={yshift=-2mm, xshift=3mm},
    boxed title style={colback=navyblue, boxrule=0pt, sharp corners, rounded corners=north},
    boxrule=1.2pt,
    sharp corners,
    rounded corners=south,
    top=4mm, bottom=3mm, left=3mm, right=3mm
}

\newtcolorbox{soluzionebox}[1][Svolgimento Dettagliato]{
    enhanced,
    breakable,
    colback=subtlegray,
    colframe=forestgreen!80!black,
    coltitle=white,
    fonttitle=\bfseries\sffamily,
    title={#1},
    attach boxed title to top left={yshift=-2mm, xshift=3mm},
    boxed title style={colback=forestgreen!80!black, boxrule=0pt, sharp corners, rounded corners=north},
    boxrule=1pt,
    sharp corners,
    rounded corners=south,
    top=4mm, bottom=3mm, left=3mm, right=3mm
}

\newtcolorbox{esamebox}[1][Consiglio per l'Esame \& Aspetti Chiave]{
    enhanced,
    breakable,
    colback=crimsonred!4!white,
    colframe=crimsonred,
    coltitle=crimsonred,
    fonttitle=\bfseries\sffamily,
    title={#1},
    attach boxed title to top left={yshift=-1.5mm, xshift=2mm},
    boxed title style={colback=white, boxrule=0.8pt, colframe=crimsonred},
    boxrule=0.8pt,
    leftrule=3.5pt,
    sharp corners,
    top=3.5mm, bottom=2.5mm, left=3mm, right=3mm
}

\lstset{
    backgroundcolor=\color{codebg},
    basicstyle=\ttfamily\small,
    breaklines=true,
    frame=single,
    rulecolor=\color{navyblue!40},
    keywordstyle=\color{navyblue}\bfseries,
    commentstyle=\color{forestgreen}\itshape,
    stringstyle=\color{warmamber},
    showstringspaces=false
}

\newcommand{\EE}{\mathbb{E}}
\newcommand{\Prob}{\mathbb{P}}
\newcommand{\Var}{\operatorname{Var}}
\newcommand{\dx}{\mathrm{d}x}

\begin{document}

\begin{center}
    {\LARGE\bfseries\color{navyblue} Università di Pisa --- Corso di Laurea in Ingegneria Gestionale}\\[0.4em]
    {\Large\bfseries Statistica I (A.A. 2021/2022) --- I Appello dell'08/06/2022}\\[0.3em]
    {\large\textsf{Soluzione Completa, Rigorosa e Commentata Passo-Passo}}
\end{center}
\vspace{0.5em}
\hrule height 1.2pt
\vspace{1.5em}

% ==============================================================================
% PROBLEMA 1
% ==============================================================================
\begin{problemabox}{Problema 1 (Testo Integrale)}
Si considerino tre variabili aleatorie di Bernoulli $X,Y,Z$ indipendenti con rispettivi parametri $p_X = 1/3$, $p_Y = 1/5$, $p_Z = 2/3$.
\begin{enumerate}[label=(\alph*)]
    \item Calcolare il valore atteso e la varianza di $2X+ Y$ e di $X\cdot Z$.
    \item Calcolare la probabilità che $X = 0$ se $X+ Y = 1$.
    \item Calcolare la funzione di massa congiunta di $X+ Y$ e $X\cdot Z$. Verificare se queste variabili aleatorie sono indipendenti.
\end{enumerate}
\end{problemabox}

\begin{soluzionebox}
\textbf{1. Momenti marginali delle variabili di Bernoulli:}
\begin{itemize}
    \item $\EE[X] = 1/3, \quad \Var(X) = (1/3)(2/3) = 2/9, \quad \EE[X^2] = 1/3$.
    \item $\EE[Y] = 1/5, \quad \Var(Y) = (1/5)(4/5) = 4/25, \quad \EE[Y^2] = 1/5$.
    \item $\EE[Z] = 2/3, \quad \Var(Z) = (2/3)(1/3) = 2/9, \quad \EE[Z^2] = 2/3$.
\end{itemize}

\textbf{2. Risoluzione Quesito (a):}
\begin{align*}
\EE[2X+Y] &= 2\EE[X] + \EE[Y] = 2\left(\frac{1}{3}\right) + \frac{1}{5} = \frac{2}{3} + \frac{1}{5} = \mathbf{\frac{13}{15} \approx 0.8667}. \\
\Var(2X+Y) &= 4\Var(X) + \Var(Y) = 4\left(\frac{2}{9}\right) + \frac{4}{25} = \frac{8}{9} + \frac{4}{25} = \frac{200 + 36}{225} = \mathbf{\frac{236}{225} \approx 1.0489}.
\end{align*}
Per il prodotto $X \cdot Z$, essendo $X$ e $Z$ indipendenti:
\begin{align*}
\EE[XZ] &= \EE[X]\EE[Z] = \left(\frac{1}{3}\right)\left(\frac{2}{3}\right) = \mathbf{\frac{2}{9} \approx 0.2222}. \\
\Var(XZ) &= \EE[(XZ)^2] - (\EE[XZ])^2 = \EE[X^2]\EE[Z^2] - \left(\frac{2}{9}\right)^2 = \left(\frac{1}{3}\right)\left(\frac{2}{3}\right) - \frac{4}{81} = \frac{2}{9} - \frac{4}{81} = \mathbf{\frac{14}{81} \approx 0.1728}.
\end{align*}

\textbf{3. Risoluzione Quesito (b):}
\[
\Prob(X=0 \mid X+Y=1) = \frac{\Prob(X=0, X+Y=1)}{\Prob(X+Y=1)} = \frac{\Prob(X=0, Y=1)}{\Prob(X=0, Y=1) + \Prob(X=1, Y=0)}.
\]
Sostituendo i valori di probabilità:
\[
\Prob(X=0 \mid X+Y=1) = \frac{\frac{2}{3} \cdot \frac{1}{5}}{\left(\frac{2}{3} \cdot \frac{1}{5}\right) + \left(\frac{1}{3} \cdot \frac{4}{5}\right)} = \frac{\frac{2}{15}}{\frac{2}{15} + \frac{4}{15}} = \frac{2}{6} = \mathbf{\frac{1}{3} \approx 0.3333}.
\]

\textbf{4. Risoluzione Quesito (c):}
I supporti sono $X+Y \in \{0, 1, 2\}$ e $XZ \in \{0, 1\}$.
Calcoliamo le probabilità congiunte:
\begin{itemize}
    \item Se $X+Y=0 \implies X=0, Y=0$, ne consegue che $XZ = 0 \cdot Z = 0$. Quindi:
    \[
    \Prob(X+Y=0, XZ=0) = \Prob(X=0, Y=0) = \frac{2}{3} \cdot \frac{4}{5} = \frac{8}{15} = \frac{24}{45}, \qquad \Prob(X+Y=0, XZ=1) = 0.
    \]
    \item Se $X+Y=1$:
    \begin{itemize}
        \item $X=0, Y=1 \implies XZ=0$: $\Prob = \frac{2}{3} \cdot \frac{1}{5} = \frac{2}{15} = \frac{6}{45}$.
        \item $X=1, Y=0 \implies XZ = Z$:
        $\Prob(XZ=0) = \frac{1}{3} \cdot \frac{4}{5} \cdot \frac{1}{3} = \frac{4}{45}$, \quad $\Prob(XZ=1) = \frac{1}{3} \cdot \frac{4}{5} \cdot \frac{2}{3} = \frac{8}{45}$.
    \end{itemize}
    Totale $X+Y=1$: $\Prob(XZ=0) = \frac{6+4}{45} = \frac{10}{45}$, \quad $\Prob(XZ=1) = \frac{8}{45}$.
    \item Se $X+Y=2 \implies X=1, Y=1 \implies XZ = Z$:
    \[
    \Prob(XZ=0) = \frac{1}{3} \cdot \frac{1}{5} \cdot \frac{1}{3} = \frac{1}{45}, \qquad \Prob(XZ=1) = \frac{1}{3} \cdot \frac{1}{5} \cdot \frac{2}{3} = \frac{2}{45}.
    \]
\end{itemize}
\begin{center}
\begin{tabular}{c|cc|c}
\toprule
$p(u, v)$ & $XZ = 0$ & $XZ = 1$ & Marginale $X+Y$ \\
\midrule
$X+Y = 0$ & $24/45$ & $0$ & $24/45 = 8/15$ \\
$X+Y = 1$ & $10/45$ & $8/45$ & $18/45 = 2/5$ \\
$X+Y = 2$ & $1/45$ & $2/45$ & $3/45 = 1/15$ \\
\midrule
Marginale $XZ$ & $35/45 = 7/9$ & $10/45 = 2/9$ & $1$ \\
\bottomrule
\end{tabular}
\end{center}
\emph{Verifica indipendenza:} $\Prob(X+Y=0, XZ=1) = 0 \neq \Prob(X+Y=0)\Prob(XZ=1) = \frac{8}{15} \cdot \frac{2}{9} = \frac{16}{135}$.
Pertanto, le variabili $X+Y$ e $XZ$ \textbf{NON sono indipendenti}.
\end{soluzionebox}

% ==============================================================================
% PROBLEMA 2
% ==============================================================================
\begin{problemabox}{Problema 2 (Testo Integrale)}
Un produttore di pneumatici dice che i suoi prodotti durano almeno il 20\% in più della media dei pneumatici in commercio ($50000\text{ km}$). Compriamo $10$ pneumatici con risultati (in km):
\[
60900, \quad 61500, \quad 64400, \quad 59500, \quad 59580, \quad 64265, \quad 63400, \quad 58500, \quad 58000, \quad 60200.
\]
Si assuma distribuzione normale $\mathcal{N}(\mu, \sigma^2)$ con media e varianza incognite.
\begin{enumerate}[label=(\alph*)]
    \item Si calcoli l'intervallo di confidenza bilaterale al 95\% per $\mu$.
    \item Si definiscano le ipotesi $H_0$ e $H_1$ per confermare il messaggio del produttore.
    \item Si testi l'ipotesi al livello di significatività del 5\%.
    \item Si calcoli approssimativamente il valore $p$.
\end{enumerate}
\end{problemabox}

\begin{soluzionebox}
\textbf{1. Statistiche descrittive ($n = 10$):}
\[
\overline{x}_{10} = \mathbf{61024.5\text{ km}}, \qquad s_{10} = \mathbf{2318.78\text{ km}}.
\]

\textbf{2. Risoluzione Quesito (a) --- Intervallo di Confidenza al 95\% ($\nu = 9$ g.d.l.):}
Il valore critico della distribuzione $t$ di Student a $9$ g.d.l. è $t_{9, \, 0.025} = 2.2622$.
\[
ME = t_{9, \, 0.025} \frac{s}{\sqrt{10}} = 2.2622 \cdot \frac{2318.78}{\sqrt{10}} = 2.2622 \cdot 733.26 \approx \mathbf{1658.75\text{ km}}.
\]
\[
\mathbf{\operatorname{IC}_{95\%}(\mu) = [61024.5 - 1658.75, \;\; 61024.5 + 1658.75] = [59365.75, \;\; 62683.25]\text{ km}}.
\]

\textbf{3. Risoluzione Quesito (b) --- Formulazione del test d'ipotesi:}
Il valore soglia dichiarato è $\mu_0 = 1.20 \times 50000 = 60000\text{ km}$.
Il messaggio del produttore viene confermato se si rigetta l'ipotesi nulla in favore dell'ipotesi alternativa:
\[
\mathbf{H_0: \mu \le 60000} \qquad \text{vs} \qquad \mathbf{H_1: \mu > 60000}.
\]

\textbf{4. Risoluzione Quesito (c) --- Esecuzione del $t$-test unilaterale destro ($\alpha = 0.05$):}
Calcoliamo la statistica test $T$:
\[
T_{\text{oss}} = \frac{\overline{x}_{10} - \mu_0}{s/\sqrt{10}} = \frac{61024.5 - 60000}{733.26} = \mathbf{1.3972}.
\]
Il valore critico unilaterale destro a $\alpha = 0.05$ con $9$ g.d.l. è $t_{9, \, 0.05} = 1.8331$.
Poiché $T_{\text{oss}} = 1.3972 < 1.8331$, il test \textbf{non rigetta $H_0$}.
\emph{Conclusione:} I dati campionari non forniscono evidenza statistica sufficiente a confermare l'affermazione pubblicitaria del produttore.

\textbf{5. Risoluzione Quesito (d) --- Calcolo del $p$-value:}
\[
p\text{-value} = \Prob(t_9 \ge 1.3972) \approx \mathbf{0.0979} \quad (\mathbf{9.79\%} \approx 10\%).
\]
\end{soluzionebox}

% ==============================================================================
% PROBLEMA 3
% ==============================================================================
\begin{problemabox}{Problema 3 (Testo Integrale)}
Si consideri una serie di variabili aleatorie i.i.d. $X_1,\dots,X_{100}$ uniformi in $[-1,5] \subset \mathbb{R}$.
\begin{enumerate}[label=(\alph*)]
    \item Calcolare la probabilità che non ci siano valori negativi tra $X_1,\dots,X_{10}$.
    \item Calcolare la probabilità che non ci siano valori negativi tra $X_1,\dots,X_{10}$ e non ci siano valori maggiori di $3$ tra $X_6,\dots,X_{13}$.
    \item Calcolare il valore atteso e la varianza di $Y = \sum_{i=1}^{100} X_i$.
    \item Calcolare approssimativamente $\Prob(Y < 170)$.
\end{enumerate}
\end{problemabox}

\begin{soluzionebox}
\textbf{1. Proprietà della singola variabile $X_i \sim \operatorname{Unif}(-1, 5)$:}
L'ampiezza dell'intervallo è $5 - (-1) = 6$.
\[
\EE[X_i] = \frac{-1 + 5}{2} = 2, \qquad \Var(X_i) = \frac{6^2}{12} = 3.
\]

\textbf{2. Risoluzione dei quesiti:}
\begin{enumerate}[label=(\alph*)]
    \item $\Prob(X_i > 0) = \frac{5 - 0}{6} = \frac{5}{6}$. Per l'indipendenza:
    \[
    \Prob(X_1 > 0, \dots, X_{10} > 0) = \left(\frac{5}{6}\right)^{10} \approx \mathbf{0.1615} \quad (\mathbf{16.15\%}).
    \]

    \item Suddividiamo le variabili in blocchi disgiunti:
    \begin{itemize}
        \item $X_1, \dots, X_5$: vincolate a $X_i \in (0, 5] \implies \Prob = 5/6$;
        \item $X_6, \dots, X_{10}$: vincolate a $X_i \in (0, 3] \implies \Prob = \frac{3 - 0}{6} = \frac{3}{6} = \frac{1}{2}$;
        \item $X_{11}, \dots, X_{13}$: vincolate a $X_i \in [-1, 3] \implies \Prob = \frac{3 - (-1)}{6} = \frac{4}{6} = \frac{2}{3}$.
    \end{itemize}
    Per reciproca indipendenza:
    \[
    \Prob(\text{Evento}) = \left(\frac{5}{6}\right)^5 \left(\frac{1}{2}\right)^5 \left(\frac{2}{3}\right)^3 = \frac{3125}{7776} \cdot \frac{1}{32} \cdot \frac{8}{27} = \frac{3125}{839808} \approx \mathbf{0.003721} \quad (\mathbf{0.3721\%}).
    \]

    \item Per la somma $Y = \sum_{i=1}^{100} X_i$:
    \[
    \EE[Y] = 100 \EE[X_i] = 100(2) = \mathbf{200}, \qquad \Var(Y) = 100 \Var(X_i) = 100(3) = \mathbf{300}.
    \]

    \item Per il TLC, $Y \approx \mathcal{N}(200, 300)$ con $\sigma_Y = \sqrt{300} = 10\sqrt{3} \approx 17.3205$:
    \[
    \Prob(Y < 170) \approx \Phi\left(\frac{170 - 200}{\sqrt{300}}\right) = \Phi\left(\frac{-30}{17.3205}\right) = \Phi(-\sqrt{3}) = \Phi(-1.7321).
    \]
    Dalle tavole: $\Phi(-1.7321) = 1 - 0.9584 = \mathbf{0.0416} \quad (\mathbf{4.16\%})$.
\end{enumerate}
\end{soluzionebox}

% ==============================================================================
% PROBLEMA 4
% ==============================================================================
\begin{problemabox}{Problema 4 (Testo Integrale)}
Sia $(X_1,\dots,X_n)$ un campione estratto da una popolazione con densità $f_\theta(x) = c_\theta \mathbf{1}_{[\theta,1]}(x)$, con $\theta \in (0,1)$.
\begin{enumerate}[label=(\alph*)]
    \item Determinare il valore di $c_\theta$ e di $\EE[X_1]$ per ogni $\theta \in [0,1)$.
    \item Determinare lo stimatore $\widehat{\theta}_n$ di massima verosimiglianza per $\theta$.
    \item Verificare se lo stimatore $\widetilde{\theta} = \max\left(-1 + \frac{2}{n}\sum_{i=1}^n X_i, \, 0\right)$ è corretto per $n=1$.
    \item Verificare se $\widetilde{\theta}$ è consistente.
\end{enumerate}
\end{problemabox}

\begin{soluzionebox}
\begin{enumerate}[label=(\alph*)]
    \item Integrando la densità:
    \[
    1 = \int_\theta^1 c_\theta \, \dx = c_\theta (1 - \theta) \implies \mathbf{c_\theta = \frac{1}{1-\theta}}.
    \]
    \[
    \EE[X_1] = \int_\theta^1 x \frac{1}{1-\theta} \, \dx = \frac{1}{1-\theta} \left[ \frac{x^2}{2} \right]_\theta^1 = \frac{1 - \theta^2}{2(1-\theta)} = \mathbf{\frac{1+\theta}{2}}.
    \]

    \item La funzione di verosimiglianza è:
    \[
    L(\theta) = \prod_{i=1}^n \frac{1}{1-\theta} \mathbf{1}_{[\theta, 1]}(x_i) = \frac{1}{(1-\theta)^n} \mathbf{1}_{\{\theta \le \min(x_1,\dots,x_n)\}}.
    \]
    Per $\theta \in (0, \min x_i]$, la funzione $\frac{1}{(1-\theta)^n}$ è strettamente crescente rispetto a $\theta$. Il massimo si trova sul confine superiore:
    \[
    \mathbf{\widehat{\theta}_{\text{MLE}} = \min(X_1, \dots, X_n) = X_{(1)}}.
    \]

    \item Per $n=1$, $\widetilde{\theta} = \max(-1 + 2X_1, 0)$.
    Calcoliamo il valore atteso per $\theta = 1/3$:
    Poiché $X_1 \ge \theta = 1/3$, per $X_1 \in [1/3, 1/2]$ si ha $-1+2X_1 \le 0 \implies \max = 0$.
    \[
    \EE[\widetilde{\theta}] = \int_{1/2}^1 (-1 + 2x) \frac{1}{1 - 1/3} \, \dx = \frac{3}{2} \left[ -x + x^2 \right]_{1/2}^1 = \frac{3}{2} \left( 0 - \left(-\frac{1}{2} + \frac{1}{4}\right) \right) = \frac{3}{2} \cdot \frac{1}{4} = \frac{3}{8} \neq \frac{1}{3}.
    \]
    Poiché $\EE[\widetilde{\theta}] = 3/8 \neq \theta = 1/3$, lo stimatore \textbf{NON è corretto}.

    \item Per la Legge Debole dei Grandi Numeri (WLLN), $\overline{X}_n \xrightarrow{\Prob} \EE[X_1] = \frac{1+\theta}{2}$.
    Per il Teorema di Continuità (Continuous Mapping Theorem), essendo la trasformazione $g(u) = \max(-1 + 2u, 0)$ continua:
    \[
    \widetilde{\theta} = g(\overline{X}_n) \xrightarrow{\Prob} g\left(\frac{1+\theta}{2}\right) = \max\left(-1 + 2\left(\frac{1+\theta}{2}\right), \, 0\right) = \max(\theta, 0) = \theta.
    \]
    Pertanto, lo stimatore $\widetilde{\theta}$ è \textbf{consistente}.
\end{enumerate}
\end{soluzionebox}

% ==============================================================================
% PROBLEMA 5
% ==============================================================================
\begin{problemabox}{Problema 5 (Testo Integrale)}
Consideriamo l'integrale
\[
\int_{-\infty}^\infty \frac{e^{-x^2/2 - 3/2}\log(7 + e^x)}{\sqrt{2\pi}(2 + \sin(x))} \, \dx.
\]
\begin{enumerate}[label=(\alph*)]
    \item Discutere l'integrazione di Monte Carlo per questo integrale.
    \item Scrivere i comandi R per approssimare l'integrale.
\end{enumerate}
\end{problemabox}

\begin{soluzionebox}
\textbf{1. Formulazione Monte Carlo tramite Gaussiana Standard:}
Riscriviamo l'integrale fattorizzando la densità normale standard $\phi(x) = \frac{1}{\sqrt{2\pi}} e^{-x^2/2}$:
\[
I = \int_{-\infty}^\infty \left[ e^{-3/2} \frac{\log(7 + e^x)}{2 + \sin(x)} \right] \frac{1}{\sqrt{2\pi}} e^{-x^2/2} \, \dx = \EE[g(X)], \qquad X \sim \mathcal{N}(0, 1),
\]
dove $g(x) = e^{-3/2} \frac{\log(7 + e^x)}{2 + \sin(x)}$. Valore di riferimento: $I \approx \mathbf{0.2644}$.

\textbf{2. Codice R:}
\begin{lstlisting}[language=R]
set.seed(42)
N <- 1000000
x_norm <- rnorm(N, mean = 0, sd = 1)
g_val <- exp(-1.5) * log(7 + exp(x_norm)) / (2 + sin(x_norm))
I_hat <- mean(g_val)
cat(sprintf("Stima Monte Carlo: %.4f\n", I_hat)) # 0.2644
\end{lstlisting}
\end{soluzionebox}

\end{document}
